\documentclass[10pt,a4paper]{amsart}
\usepackage[margin=19mm]{geometry}
\usepackage{amsmath,amssymb,mathtools}
\usepackage[scr=rsfs,cal=euler]{mathalfa}
\usepackage{xfrac}
\usepackage[unicode,hidelinks,bookmarks=false]{hyperref}
\newcommand{\ie}{i.e.\ }
\newcommand{\cf}{cf.\ }
\newcommand{\resp}{respectively\ }
\newcommand{\Subsection}{\subsection}
\providecommand{\nobreakdash}{\nobreak\hbox{-}}
\setlength{\emergencystretch}{2em}
\multlinegap0pt
\usepackage{mathtools}
\usepackage{amssymb}
\usepackage{enumerate}
\usepackage{xcolor}
\newtheorem{lemma}{Lemma}
\newcommand{\R}{\mathbb{R}}
\title{Relaxation dynamics of the Inertial Winfree model}
\author{Caiman Moreno-Earle, Seung-Yeon Ryoo, Grace To}
\date{Source: arXiv:2605.01695v1 (CC BY 4.0)}
\begin{document}
\maketitle

\noindent\emph{Source and licence.} Relaxation dynamics of the Inertial Winfree model. Version arXiv:2605.01695v1 (CC BY 4.0), distributed by arXiv under the Creative Commons Attribution 4.0 International licence, http://creativecommons.org/licenses/by/4.0/, as recorded in the arXiv licence metadata for this exact version and shown on the abstract page https://arxiv.org/abs/2605.01695v1. Full licence notice: \texttt{licenses/arxiv-2605.01695-CC-BY-4.0.txt}.\par\smallskip

\noindent\textit{Editorial scope.} Counted unit: the article's lemma lem:single (source0.tex lines 765--779). The article's complete original proof of it is delivered with it (source0.tex lines 787--847, one \texttt{proof} environment of 61 lines, which the article places away from the statement). No mathematics is written, repaired or re-derived: the statement and the proof are the article's own lines, in the article's own notation, and no retained line is substituted or re-flowed.  Retained uncounted as the source-local context the proof needs: lines 117--122; lines 439--451.  Nothing was omitted from any retained span, so the package carries no omission record.  No retained line contains a citation, so the wrapper carries no bibliography and no bibliography entry.  No retained line contains a figure environment, an image-inclusion command or drawing code, so no image is delivered and the package declares no asset dependency.  Editorial formatting only, in addition: an AMS \texttt{amsart} preamble that restates the article's own declarations and macros used by the retained lines, namely the environments \texttt{lemma} on separate counters, together with the standard packages those lines need.  The retained environments print in this document as Lemma 1 in order of appearance, whereas the article prints the same items with the numbering of its own full document.  The complete native source of the article as distributed in the arXiv e-print for this exact version is bundled unchanged as \texttt{sources/199654/source0.tex}.
\begin{equation}\label{Winfree}
    \begin{cases}
        m\ddot\theta_i + \dot{\theta}_i=\nu_i-\dfrac\kappa N \sum_{j=1}^N (1+\cos\theta_j(t))\sin\theta_i(t), & t>0, \\[2pt]
        (\theta_i,\dot\theta_i)\big\rvert_{t=0^+} =(\theta_i^0,\omega_i^0), & i=1,\dots,N.
    \end{cases}
\end{equation}

\begin{lemma}[Approximating second-order velocity by first-order]\label{lem:orderlowering}
For every solution to~\eqref{Winfree}, every $i\in[N]$, and every $t\ge 0$,
\begin{multline*}
\bigl|\dot\theta_i(t)-\omega_i^0 e^{-t/m}-\nu_i(1-e^{-t/m})+\kappa R(t)\sin\theta_i(t)(1-e^{-t/m})\bigr|\\
\le 2\kappa\,\|\Omega^0\|_\infty t e^{-t/m}(1-e^{-t/m})+2m\kappa M_F(1-e^{-t/m})^{3}.
\end{multline*}

If $\mathcal{B}\subset [N]$ and $i\in \mathcal{B}$, then
\begin{multline*}
	\bigl|\dot\theta_i(t)-\omega_i^0 e^{-t/m}-\nu_i(1-e^{-t/m})+\kappa R_\mathcal{B}(t)\sin\theta_i(t)(1-e^{-t/m})\bigr|\\
	\le 2\kappa (1-e^{-t/m})\left(\frac{|\mathcal{B}|}{N} \left(\|\Omega_\mathcal{B}^0\|_\infty te^{-t/m}+m M_{F,\mathcal{B}}(1-e^{-t/m})^2\right)+\frac{N-|\mathcal{B}|}{N} \right).
\end{multline*}
\end{lemma}

\begin{lemma}[A priori partial trapping]\label{lem:single}
Let $\mathcal{B}\subset [N]$, $\eta>0$, $\rho>0$, and $T\in (\eta m,\infty]$, and suppose $R_\mathcal{B}(t)\ge\rho$ for
$t\in[\eta m,T)$. Let $x\in (-1,1)$ satisfy
\begin{equation}\label{eq:xbound}
1-x^2\ge\left(\frac{\|\Omega_\mathcal{B}^0\|_\infty e^{-\eta}+\|\mathcal V_\mathcal{B}\|_\infty+2\kappa\left(\frac{|\mathcal{B}|}{N}m\left(\|\Omega_\mathcal{B}^0\|_\infty(\eta\vee 1)e^{-(\eta\vee 1)}+ M_{F,\mathcal{B}}\right)+\frac{N-|\mathcal{B}|}{N}\right)}{\kappa\rho(1-e^{-\eta})}\right)^{\!2},
\end{equation}
where $z\vee w\coloneqq \max\{z,w\}$.
Let $y\in[-|x|,|x|]$.
\begin{enumerate}
\item If $i\in \mathcal{B}$ is such that $\cos\theta_i(t_0)\ge y$ for
some $t_0\in[\eta m,T)$, then $\cos\theta_i(t)\ge y$ for all
$t\in[t_0,T)$. In particular, $\mathrm{osc}_{[\eta m,T)} \theta_i<2\pi$, where $\mathrm{osc}_{I}f(t)\coloneqq \sup_{t\in I}f(t)-\inf_{t\in I} f(t)$ for an interval $I$ and a function $f:I\to \R$.
\item For all $i\in \mathcal{B}$, we have $\mathrm{osc}_{[\eta m,T)} \theta_i<4\pi$.
\end{enumerate}
\end{lemma}

\begin{proof}[Proof of Lemma \ref{lem:single}]
	\,
\begin{enumerate}
\item 
If $\cos\theta_i(t)>y$ for all $t\in [t_0,T)$, then there is nothing to prove. Otherwise, let $t_1\in[t_0,T)$ be the infimal time  such that
$\cos\theta_i(t_1)\le y$; this necessitates $\cos\theta_i(t_1)=y$ by continuity. We claim that $\frac{d}{dt}\cos\theta_i(t_1)>0$. Indeed, at $t=t_1$, $|\sin\theta_i(t_1)|=\sqrt{1-y^2}\ge\sqrt{1-x^2}$. Using
Lemma~\ref{lem:orderlowering} at $t_1$,
\begin{align*}
\bigl.\tfrac{d}{dt}\cos\theta_i\bigr|_{t=t_1}
&=-\sin\theta_i(t_1)\dot\theta_i(t_1)\\
&\ge \kappa R_\mathcal{B}(t_1)\sin^2\theta_i(t_1)(1-e^{-t_1/m})\\
&\quad-|\sin\theta_i(t_1)|\bigl(\|\Omega_\mathcal{B}^0\|_\infty e^{-t_1/m}+\|\mathcal V_\mathcal{B}\|_\infty(1-e^{-t_1/m})\bigr)\\
&\quad-|\sin\theta_i(t_1)|2\kappa (1-e^{-t_1/m})\left(\frac{|\mathcal{B}|}{N}\bigl(\|\Omega_\mathcal{B}^0\|_\infty t_1 e^{-t_1/m}+m M_{F,\mathcal{B}}(1-e^{-t_1/m})^2\bigr)+\frac{N-|\mathcal{B}|}{N}\right).
\end{align*}

Dividing by $|\sin\theta_i(t_1)|=\sqrt{1-y^2}(\ge\sqrt{1-x^2}>0)$ and using
$R_\mathcal{B}(t_1)\ge\rho$, $t_1\ge\eta m$,
\begin{align*}
\frac{1}{\sqrt{1-y^2}}\bigl.\tfrac{d}{dt}\cos\theta_i\bigr|_{t=t_1}
&\ge \kappa\rho\sqrt{1-y^2}(1-e^{-\eta})-\|\Omega_\mathcal{B}^0\|_\infty e^{-\eta}-\|\mathcal V_\mathcal{B}\|_\infty\\
&\quad-2\kappa\left(\frac{|\mathcal{B}|}{N}m\left(\|\Omega_\mathcal{B}^0\|_\infty(\eta\vee 1)e^{-(\eta\vee 1)}+ M_{F,\mathcal{B}}\right)+\frac{N-|\mathcal{B}|}{N}\right)\\
&>0,
\end{align*}
where the first inequality uses the fact applied to $s=t_1/m$ that the function $s\mapsto se^{-s}$ is unimodal with maximum at $s=1$, so that $s\ge \eta$ implies that $se^{-s}\le (\eta\vee 1)e^{-(\eta\vee 1)}$, and
the second inequality is \eqref{eq:xbound} combined with $1-y^2\ge 1-x^2$. Thus
$\frac{d}{dt}\cos\theta_i(t_1)>0$.

Since $\cos\theta_i(t_0)\ge y$ yet $t_1$ is the first time such that $\cos\theta_i(t_1)\le y$, the inequality $\frac{d}{dt}\cos\theta_i(t_1)>0$ implies that $t_1=t_0$ (since otherwise $t_1>t_0$ and there are times slightly smaller than $t_1$ at which $\cos \theta_i<y$, violating the minimality of $t_1$) So, it must be that $\cos\theta_i(t_0)=y$ and $\cos\theta_i(t)>y$ for $t\in [t_0,t_0+\varepsilon)\subset [t_0,T)$ for some small $\varepsilon$.

If there were another time $t_2\in [t_0+\varepsilon,T)$ for which $\cos\theta_i(t_2)\le y$, then by making $t_2$ the earliest such time, we must have $\cos\theta_i(t_2)=y$, and the same computation as above shows that $\frac{d}{dt}\cos\theta_i(t_2)>0$, giving a contradiction (as there would be times slightly smaller than $t_2$ at which $\cos\theta_i<y$, violating the minimality of $t_2$).

By this exit-time argument, we have $\cos\theta_i(t)\ge y$ for all $t\in[t_0,T)$ (and in fact, the stronger statement that $\cos\theta_i(t)>y$ for $t\in (t_0,T)$).

The latter statement follows from $y>-1$, and the continuity of $\theta_i$ in $t$. More explicitly, $\theta_i|_{[t_1,T)}$ is confined to a single component of $\{\cos\theta\ge -|x|\}$, a closed
interval of length $2(\pi-\arccos|x|)<2\pi$; hence
$\mathrm{osc}_{[t_1,T)}\,\theta_i<2\pi$. 
\item Fix $i\in\mathcal{B}$. Define
$t^{*}:=\inf\{t\in[\eta m,T):\cos\theta_i(t)\ge -|x|\}$
(set $t^{*}=T$ if the set is empty).

If $t^{*}=T$, then $\cos\theta_i(t)<-|x|$ for every $t\in[\eta m,T)$;
since $|x|<1$, the continuous lift $\theta_i$ is confined to a single
connected component of $\{\theta\in\R:\cos\theta<-|x|\}$, an open
interval of length $2\arccos|x|<2\pi$.

If $t^{*}<T$, then by continuity $\cos\theta_i(t^{*})\ge -|x|$, and
part~(1) with $y=-|x|$ gives $\cos\theta_i(t)\ge -|x|$ for all
$t\in[t^{*},T)$. The lift $\theta_i|_{[t^{*},T)}$ is therefore
confined to a single component of $\{\cos\theta\ge -|x|\}$, a closed
interval of length $2(\pi-\arccos|x|)<2\pi$; hence
$\mathrm{osc}_{[t^{*},T)}\,\theta_i<2\pi$. On $[\eta m,t^{*})$ we
have $\cos\theta_i<-|x|$, so by the same connected-component argument
$\mathrm{osc}_{[\eta m,t^{*}]}\,\theta_i\le 2\arccos|x|<2\pi$.
Sub-additivity of oscillation gives
\[
\mathrm{osc}_{[t_1,T)}\,\theta_i\le \mathrm{osc}_{[\eta m,t^{*}]}\,\theta_i+\mathrm{osc}_{[t^{*},T)}\,\theta_i
<2\pi+2\pi=4\pi.\qedhere
\]

\end{enumerate}
\end{proof}

\end{document}
